\section{Filtros de tiempo discreto}

Un filtro digital es un algoritmo que, ejecutado en una computadora,
convierte una secuencia de n\'{u}meros $x[n]$ de entrada, en otra
secuencia $y[n]$ de salida. Adem\'{a}s de filtrar
ciertas bandas de frecuencia, puede emplearse para otras funciones
tales como integraci\'{o}n, diferenciaci\'{o}n, y estimaci\'{o}n.

La ecuaci\'{o}n de diferencias que relaciona entrada y salida puede
ser expresado en el dominio del tiempo discreto como una suma de la
forma
\begin{equation}
y[n] = \sum_{i = 0}^{k} a_{i} \, x [n - i] - \sum_{i=0}^{k} b_{i} \, y[n-i],
\end{equation}
o en el dominio $z$ como
\begin{equation}
G(z) = \frac{\sum_{i = 0}^{k} a_{i} \, z^{-i}}{1 + \sum_{i=0}^{k} b_{i} \, z^{-i}}.
\end{equation}

Por lo tanto, el dise\~{n}o de un filtro digital implica la
determinaci\'{o}n de los coeficientes $a_{i}$ y $b_{i}$.

\subsection{Filtro recursivo o de respuesta al impulso infinita \textbf{IIR}}

Sea el sistema de tiempo discreto con la ec. de diferencias 
\begin{eqnarray}
        y\left[ n\right] -ay\left[ n-1\right]   &=&     x\left[ n\right] , \\
        \left| a\right|                         &<&     1.
\end{eqnarray}
Como $y\left[ n\right] =H\left( e^{j\omega }\right) e^{j\omega n}$ se obtiene 
\begin{equation}
        H\left( e^{j\omega }\right) e^{j\omega n}-aH\left( e^{j\omega }\right) e^{j\omega n-1} = 
                e^{j\omega n}.
\end{equation}
Agrupando y simplificando 
\begin{equation}
        \left( 1-ae^{-j\omega }\right) H\left( e^{j\omega }\right) =1.
\end{equation}
Por lo tanto 
\begin{equation}
        H\left( e^{j\omega }\right) =\frac{1}{1-ae^{-j\omega }}.
\end{equation}
La respuesta al impulso es 
\begin{equation}
        h[n]=a^nu\left[ n\right] .
\end{equation}
La respuesta al escal\'{o}n es 
\begin{equation}
        s\left[ n\right] =\frac{1-a^{n+1}}{1-a}u\left[ n\right] .
\end{equation}

\subsection{Filtros de respuesta al impulso finita \textbf{FIR}}

La forma general de una ecuaci\'{o}n de diferencias no recursiva es 
\begin{equation}
        y\left[ n\right] =\sum_{k=-N}^{M}b_{k}x\left[ n-k\right].
\end{equation}

Por ejemplo, un filtro promediador 
\begin{equation}
        y\left[ n\right] =\frac{1}{M+N+1}\sum_{k=-N}^{M}x\left[ n-k\right] .
\end{equation}

\begin{equation}
        h\left[ n\right] =      \left\{ 
                                \begin{array}{ll}
                                        \frac{1}{N+M+1} & -N\leq n\leq M, \\ 
                                        0 & n<-N, \\ 
                                        0 & M<n.
                                \end{array}
                                \right.
\end{equation}
Recordando que 
\begin{equation}
        H\left( e^{j\omega }\right) =\sum_{k=-\infty }^{\infty }h\left[ k\right]
e^{-j\omega k},
\end{equation}
resulta 
\begin{eqnarray}
       H\left( e^{j\omega }\right) &=&\frac{1}{M+N+1}\sum_{k=-\infty }^{\infty}e^{-j\omega k}, \\
       H\left( e^{j\omega }\right) &=&
   \frac{1}{M+N+1}e^{j\omega \frac{\left(N-M\right) }{2}}
   \frac{\sin \left( \frac{\omega \left( M+N+1\right) }{2}\right) }{\sin \left( \frac{\omega }{2}\right) }.
\end{eqnarray}

Otro filtro, pasaaltos 
\begin{eqnarray}
      y\left[ n\right]                &=&\frac{x\left[ n\right] -x\left[ n-1\right] }{2}, \\
      H\left( e^{j\omega }\right)     &=&\frac{1}{2}\left( 1-e^{-j\omega }\right) =
                                             je^{\frac{j\omega }{2}}\sin \left( \frac{\omega }{2}\right) .
\end{eqnarray}

En general, supongamos una funci\'{o}n de transferencia del tipo
\begin{equation}
        H\left( z\right) =\sum_{k=0}^{M}h\left[ k\right] z^{-k},
\end{equation}
correspondiente a la ecuaci\'{o}n de diferencias 
\begin{equation}
        y\left[ n\right] =\sum_{k=0}^{M}h\left[ k\right] x\left[ n-k\right] .
\end{equation}
Este tipo de filtro es conocido como $FIR$, o de \textbf{R}espuesta al\textbf{I}mpulso\textbf{F}inita.

Observemos que :
\begin{enumerate}
      \item  Los coeficientes de $H\left( z\right) $ son los valores de la respuesta al impulso del filtro.
      \item  La respuesta al impulso tiene una duraci\'{o}n de $M+1$ muestras.
\end{enumerate}

\subsection{Dise\~{n}o de filtros de tiempo discreto con respuesta al impulso infinita (IIR) }

\begin{enumerate}
        \item   Para dise\~{n}ar un filtro debemos definir las especificaciones
        \begin{enumerate}
        
                \item   Ancho y ubicaci\'{o}n de la banda de paso;
                
                \item   Ancho y ubicaci\'{o}n de la banda de transici\'{o}n;
                
                \item   Ubicaci\'{o}n de la banda de rechazo;
                
                \item   Atenuaci\'{o}n en la banda de rechazo.

                \item   Orden del filtro $n$.
                
        \end{enumerate}

        \item   Es necesario adem\'{a}s, que los filtros de tiempo discreto que dise\~{n}emos
          sean semejantes a los filtros de tiempo continuo que hemos visto en el cap\'{\i}tulo anterior.
          Es posible lograr este objetivo mediante una transformaci\'{o}n de variables que nos permite
          obtener versiones discretizadas de filtros an\'{a}logicos. Dichas transformaciones se obtienen
          mediante aproximaciones de la derivada y de la integral.

        \item   En Matlab disponemos de las siguientes funciones para dise\~{n}ar filtros IIR
        \begin{enumerate}

                \item   butter  - filtros Butterworth;

                \item   cheby1  - filtros Chebyshev tipo I (ripple u ondulacion en la banda de  paso);

                \item   cheby2  - filtros Chebyshev tipo II (ripple u ondulacion en la banda de  rechazo);

                \item   ellip   - filtros el\'{\i}pticos (fase maximamente plana);
       \end{enumerate}
        Debemos tener cuidado pues con la misma  funci\'{o}n se dise\~{n}an filtros an\'{a}logicos y 
        discretos.
        
        \item   Un par\'{a}metro necesario para todos los filtros es $n$, que determina la cantidad de
        polos del filtro. Cuando m\'{a}s alto sea, el filtro tendr\'{a} mayor atenuaci\'{o}n en la banda
        de rechazo, mayor pendiente en la zona de transici\'{o}n, y sera m\'{a}s complicado. Por lo tanto
        es necesario que $N$ sea lo m\'{a}s bajo posible. 
        Para poder calcularlo seg\'{u}n nuestras necesidades podemos usar las funciones
        \begin{enumerate}
        
            \item   buttord     - estimaci\'{o}n orden Butterworth;
        
            \item   cheb1ord    - estimaci\'{o}n orden Chebyshev tipo I;
        
            \item   cheb2ord    - estimaci\'{o}n orden Chebyshev tipo II;
        
            \item   ellipord    - estimaci\'{o}n orden filtro el\'{\i}ptico.
        
        \end{enumerate}

        \item   Para especificar las frecuencias de paso y de corte de los distintos filtros, debemos fijar
        primero la frecuencia de muestreo $\omega_{m}$ que no se  mide en Hertz sino en muestras por
        segundo. Por lo tanto, el per\'{\i}odo de muestreo es $T_{m} = \frac{1}{\omega_{m}}$.
        
        \item   Las frecuencias de paso y de corte de los distintos filtros se especifican como una
        fracci\'{o}n de la mitad de la frecuencia de muestreo. Por lo tanto, debemos usar n\'{u}meros
        $0 < \omega < 1$, donde $1$ representa una frecuencia  igual  a la mitad de la  frecuencia  de
        muestreo.

        \item Las funciones de transferencia que obtenemos usando estas funciones de Matlab tienen la forma
        \begin{equation}
                H(z) = \frac{b(z)}{a(z)} =
                \frac{b_{0} z^{N} + b_{1} z^{N-1}+...+ b_{N}}{a_{0} z^{N} + a_{1} z^{N-1} +...+ a_{n}}.
        \end{equation}
        donde $b(z)$ y $a(z)$ son polinomios en $z$.
        Las funciones devuelven dos vectores $b$ y $a$ de longitud $N+1$  ordenados de la siguiente forma
        \begin{eqnarray}
                b &=& [b_{0}, b_{1}, b_{2},...,b_{N}],  \\
                a &=& [a_{0}, a_{1}, a_{2},...,a_{N}].
        \end{eqnarray}
       
\end{enumerate}

\section{Funciones Matlab usadas para el dise\~{n}o de filtros}

\begin{enumerate}
        \item   Para dise\~{n}ar  filtros Butterworth  usaremos la   funci\'{o}n \textbf{butter}.
        La ayuda traducida y explicada que corresponde a esta funci\'{o}n es la siguiente
        \begin{enumerate}
                \item   butter dise\~{n}o de filtros \textbf{digitales y anal\'{o}gicos} Butterworth
\begin{verbatim}
butter Butterworth digital and analogic filter design.
\end{verbatim}  
                
                \item   Filtros Butterworth digitales pasabajos de orden N y frecuencia de corte Wn,
\begin{verbatim}
[B,A] = butter(N,Wn) 
\end{verbatim}
                Obtenemos un filtro de orden $N$ y funci\'{o}n de transferencia $\frac{B(z)}{A(z)}$. Los
                coeficientes est\'{a}n almacenados en los vectores $B$ (numerador) y $A$ (denominador), 
                ordenados en potencias descendentes de z. Las frecuencias de corte $Wn$ debe  estar normalizada
\begin{verbatim}
0.0 < Wn < 1.0
\end{verbatim}
                donde $1.0$ representa la mitad de la frecuencia de muestreo

                \item   Filtros Butterworth digitales pasabanda de orden N y banda de paso
                $W1 < \omega < W2$.
                Tenemos que escribir
\begin{verbatim}
[B,A] = butter(N,[W1 W2])                       
\end{verbatim}
              para dise\~{n}ar un filtro pasabanda, donde W1 y W2 son las frecuencias inferior y  superior
              de la banda de paso normalizadas.
        
                \item   Filtros Butterworth digitales pasaalto de orden N y frecuencia de corte Wn
\begin{verbatim}
[B,A] = butter(N,Wn,'high')
\end{verbatim}
                
                \item   Filtros Butterworth digitales suprimebanda de orden N y frecuencia de corte Wn
                Tenemos que escribir
\begin{verbatim}
[B,A] = butter(N,[W1 W2],'stop')                        
\end{verbatim}
                para dise\~{n}ar un filtro suprimebanda, donde W1 y W2 son las frecuencias inferior y
                superior de la banda de rechazo.

                \item   Ceros y polos de los filtros Butterworth
                Si queremos obtener los ceros y polos de los filtros Butterworth, debemos usar
                tres variables a la izquierda del signo $=$. Por ejemplo
\begin{verbatim}
[Z,P,K] = butter(N,Wn,'stop')
\end{verbatim}
              son los ceros, polos y ganancia de un filtro suprimebanda de orden N y banda rechazada entre
              W1 y W2.
        
                \item   Filtros anal\'{o}gicos
\begin{verbatim}
 butter(N,Wn,'s')
 butter(N,Wn,'high','s')
 butter(N,Wn,'stop','s')
\end{verbatim}
                son las \'{o}rdenes que generan filtros de tiempo continuo.
        \end{enumerate}
        
        \item   Para dise\'{n}ar  filtros Chebyshev tipo I usaremos la   funci\'{o}n \textbf{cheby1}.
        La ayuda traducida y explicada que corresponde a esta funci\'{o}n es la siguiente
        \begin{enumerate}
                \item cheby1 Dise\~{n}o de filtros digitales y anal\'{o}gicos de tipo I
\begin{verbatim}
cheby1 
  Chebyshev type I digital and analog filter design.
\end{verbatim}
                
                \item   Filtros Chebyshev Tipo I digitales pasabajo de orden N, con tolerancia R
                        (en decibeles) en la banda de paso  y frecuencia de corte Wn. 
                        Usar $R = 0.5$ como valor inicial.
\begin{verbatim}
[B,A] = cheby1(N,R,Wn)
\end{verbatim}
                
                
                \item   Filtros Chebyshev Tipo I digitales pasabanda de orden N, con tolerancia R
                        (en decibeles) en la banda de paso  y frecuencias de corte W1 (inferior)
                        y W2  (superior). Empezar con R = 0.5 si se duda que valor elegir.
\begin{verbatim}
[B,A] = cheby1(N,R,[W1 W2],'pass')
\end{verbatim}
                
                \item   Filtros Chebyshev Tipo I digitales pasaalto de orden N, con tolerancia R
                        (en decibeles) en la banda de paso y  frecuencia de corte Wn
\begin{verbatim}
[B,A] = cheby1(N,R,Wn,'high')
\end{verbatim}
                
                \item   Filtros Chebyshev Tipo I digitales suprimebanda de orden N, con tolerancia  R
                        (en decibeles) en la banda de paso y frecuencias de corte W1 (m\'{\i}nima)y W2
                        (m\'{a}xima) de la banda de rechazo
\begin{verbatim}
[B,A] = cheby1(N,R,[W1 W2],'stop')
\end{verbatim}

                \item   Ceros y polos de los filtros Chebyshev tipo I
\begin{verbatim}
[Z,P,K] = Cheby1(...)
\end{verbatim}
              Es decir, si queremos obtener los ceros y polos de los filtros Chebyshev tipo I, debemos usar
              tres variables a la izquierda del signo $=$. Por ejemplo
\begin{verbatim}
[Z,P,K] = cheby1(N,Wn,'stop')
\end{verbatim}
              son los ceros, polos y ganancia de un filtro suprimebanda de orden N y banda rechazada entre
              W1 y W2.

               \item   Filtros anal\'{o}gicos
\begin{verbatim}
cheby1(N,Wn,'s'),
cheby1(N,Wn,'high','s')
cheby1(N,Wn,'stop','s')
\end{verbatim}
        \end{enumerate}

        \item   Para dise\'{n}ar  filtros Chebyshev tipo II usaremos la   funci\'{o}n \textbf{cheby2}.
        La ayuda traducida y explicada que corresponde a esta funci\'{o}n es la siguiente
                \begin{enumerate}
                \item cheby2 Dise\~{n}o de filtros digitales y anal\'{o}gicos de tipo II
\begin{verbatim}
cheby2 
   Chebyshev type II digital and analog filter design.
\end{verbatim}
                
                \item   Filtros Chebyshev Tipo II digitales pasabajo de orden N, con atenuaci\'{o}n
                        m\'{\i}nima R (en decibeles) en la banda de rechazo  y frecuencia de corte Wn.
                        Empezar con R = 20 si se duda que valor elegir.
\begin{verbatim}
[B,A] = cheby2(N,R,Wn)
\end{verbatim}
                
                \item   Filtros Chebyshev Tipo II digitales pasabanda de orden N, con atenuaci\'{o}n
                        m\'{\i}nima R (en decibeles) en la banda de rechazo y frecuencias de corte W1
                        (inferior) y W2  (superior).
\begin{verbatim}
[B,A] = cheby2(N,R,[W1 W2],'stop')
\end{verbatim}
                
                \item   Filtros Chebyshev Tipo II digitales pasaalto de orden N, con atenuaci\'{o}n
                        m\'{\i}nima R (en decibeles) en la banda de rechazo y  frecuencia de corte Wn
\begin{verbatim}
[B,A] = cheby2(N,R,Wn,'high')
\end{verbatim}
                
                \item   Filtros Chebyshev Tipo II digitales suprimebanda de orden N, con atenuaci\'{o}n
                        m\'{\i}nima R (en decibeles) en la banda de rechazo y frecuencias de corte W1
                        (m\'{\i}nima)y W2 (m\'{a}xima) de la banda de rechazo
\begin{verbatim}
[B,A] = cheby2(N,R,[W1 W2],'stop')
\end{verbatim}

                \item   Ceros y polos de los filtros Chebyshev tipo II
\begin{verbatim}
[Z,P,K] = Cheby2(...)
\end{verbatim}
             Es decir, si queremos obtener los ceros y polos de los filtros Chebyshev tipo II, debemos usar
             tres variables a la izquierda del signo $=$. Por ejemplo
\begin{verbatim}
[Z,P,K] = cheby2(N,Wn,'stop')
\end{verbatim}
             son los ceros, polos y ganancia de un filtro suprimebanda de orden N y banda rechazada entre
             W1 y W2.

                \item   Filtros anal\'{o}gicos
\begin{verbatim}
cheby2(N,Wn,'s'),
cheby2(N,Wn,'high','s')
cheby2(N,Wn,'stop','s')
\end{verbatim}
        \end{enumerate}

        \item   Para dise\'{n}ar  filtros Chebyshev tipo II usaremos la   funci\'{o}n \textbf{cheby2}.
        La ayuda traducida y explicada que corresponde a esta funci\'{o}n es la siguiente
        \begin{enumerate}
              \item   Dise\~{n}o de filtros el\'{\i}pticos  digitales y anal\'{o}gicos
        
              \item   Filtros el\'{\i}pticos  digitales pasabajo de orden N, con  tolerancia Rp en la
                      banda de paso y atenuaci\'{o}n m\'{\i}mina Rs en la banda de rechazo,  con
                      frecuencia de corte Wn. Comenzar con $Rp = 0.5$ y $Rs = 20$ si no se sabe que elegir.
\begin{verbatim}
[B,A] = ellip(N,Rp,Rs,Wn)
\end{verbatim}
                
                \item   Filtros el\'{\i}pticos  digitales pasabanda de orden N, con  tolerancia Rp en la
                        banda de paso y atenuaci\'{o}n m\'{\i}mina Rs en la banda de rechazo,  con
                        frecuencias de corte W1 (m\'{\i}nima) y W2 (m\'{a}xima) de la banda de paso.
                        Comenzar con $Rp = 0.5$ y $Rs = 20$ si no se sabe que elegir.
                        Debemos usar el siguiente comando
\begin{verbatim}
[B,A] = ellip(N,Rp,Rs,[W1  W2])
\end{verbatim}
                        
                \item   Filtros el\'{\i}pticos  digitales pasaalto de orden N, con  tolerancia Rp en la
                        banda de paso y atenuaci\'{o}n m\'{\i}mina Rs en la banda de rechazo,  con
                        frecuencia de corte Wn.
\begin{verbatim}
[B,A] = ellip(N,Rp,Rs,Wn,'high')
\end{verbatim}

                \item   Ceros y polos de los filtros el\'{\i}pticos
\begin{verbatim}
[Z,P,K] = ellip(...)
\end{verbatim}
                Es decir, si queremos obtener los ceros y polos de los filtros el\'{\i}pticos, debemos usar
                tres variables a la izquierda del signo $=$. Por ejemplo
\begin{verbatim}
[Z,P,K] = ellip(N,Rp,Rs,Wn,'stop')
\end{verbatim}
               son los ceros, polos y ganancia de un filtro suprimebanda de orden N y banda rechazada entre
               W1 y W2.

                \item   Filtros anal\'{o}gicos
\begin{verbatim}
ellip(N,Rp,Rs,Wn,'s'),
ellip(N,Rp,Rs,Wn,'high','s')
ellip(N,Rp,Rs,Wn,'stop','s')
\end{verbatim}
        \end{enumerate}

\end{enumerate}

%\subsection{Ejemplo de dise\~{n}o de filtros de tiempo discreto}
%
%\subsubsection{Filtro Butterworth pasabajo}
%
%\begin{verbatim}
%% filtro Butterworth pasabajo de tiempo discreto
%
%Wm = 8500;  % frecuencia de muestreo
%Wmm = Wm/2; % mitad de la frecuencia de muestreo
%W1=1500;    % frecuencia superior banda de paso                     
%Wp=W1/Wmm;  % frecuencia superior banda de paso normalizada 
%rp=3;       % ripple banda de paso en dB
%W2=1.1*W1;  % frecuencia inferior banda de rechazo 
%Ws=W2/Wmm;  % frecuencia inferior banda de rechazo normalizada
%rs=20;      % atenuacion banda de rechazo en dB 
%
%% usamos buttord para dimensionar el filtro
%[N,Wn]=buttord(Wp,Ws,rp,rs);
%
%% calculamos el filtro
%[b,a]=butter(N,Wn);
%
%% graficamos respuesta en frecuencia
%dbode(b,a,1/Wm);
%\end{verbatim}
%
%\subsubsection{Filtro Butterworth pasabanda}
%
%\begin{verbatim}
%%   FILTRO Butterworth pasabanda
%Wm  = 11000;    % frecuencia de muestreo
%Wmm = Wm/2;     % mitad de la frecuencia de muestreo
%W1  = 2800;     % frecuencia minima banda de paso
%W2  = 4000;     % frecuencia maxima banda de paso
%W0 = 0.9*W1;    % frecuencia maxima banda de rechazo inferior
%W3 = 1.1*W2;    % frecuencia minima banda de rechazo superior
%rp = 3;         % ripple banda de paso
%rs = 20;        % atenuacion banda de rechazo
%
%% usamos buttord para dimensionar el filtro
%[N,Wn]=buttord([W1/Fmm W2/Fmm],[W0/Fmm W3/Fmm],rp,rs)
%
%% calculamos el filtro
%[b,a]=butter(N,Wn);
%
%% graficamos respuesta en frecuencia
%dbode(b,a,1/Wm);
%\end{verbatim}
%
%\subsubsection{Filtro eliptico pasabanda}
%
%\begin{verbatim}
%Wm = 100;               % frecuencia de muestreo
%t  = (1:100)/Wm;        % vector tiempo    
%s1 = sin(2*pi*t*5);     % componente de la se~nal a filtrar
%s2 = sin(2*pi*t*15);    % componente de la se~nal a filtrar
%s3 = sin(2*pi*t*30);    % componente de la se~nal a filtrar
%s  = s1 + s2 + s3;      % se~nal a filtrar
%plot(t,s);              % grafico de la se~nal a filtrar
%
%% Para dise~nar un filtro para mantener la sinusoide de 
%% 15 Hz y eliminar las sinusoides de 5 y 30 Hz, vamos a 
%% usar un filtro eliptico con una pasabanda de 10 a 20 Hz
%[b,a]=ellip(4,0.1,40,[10 20]*2/Wm); 
%%^ dise~no de filtro de tiempo discreto
%[H,w]=freqz(b,a,512);               
%%^ calculo de la respuesta en frecuencia
%plot(w*Wm/(2*pi),abs(H));           
%%^ grafico de la respuesta en frecuencia
%sf = filter(b,a,s);                 
%%^ filtrado
%plot(t,sf);                         
%%^ grafico de la señal filtrada
%xlabel('Tiempo (segundos)');        
%%^ rotulo x
%ylabel('Señal filtrada');           
%%^ rotulo y
%axis([0,1,-1,1]);                   
%%^ marco del grafico
%
%S=fft(s,512);                       
%%^ transformada rapida de Fourier
%SF=fft(sf,512);                     
%%^ transformada rapida de Fourier
%w=(0:255)/256*(Wm/2);               
%%^ vector de frecuencias para graficar el espectro
%                                    
%% Observe que los picos de 5 y 30 Hz han sido eliminados
%plot(w,abs[S(1:256)' SF(1:256)']));
%xlabel('Frecuencia (Hz)');
%ylabel('Magnitud de la transformada de Fourier'); 
%\end{verbatim}
%
%\subsubsection{Filtro eliptico suprimebanda}
%
%\begin{verbatim}
%Wm=11000;
%Wmm=Wm/2;
%W1=2800;
%W2=4000;
%W0=1.1*W1;
%W3=0.9*W2;
%rp=3;
%rs=40;
%
%[N,Wn]=ellipord([W1/Wmm W2/Wmm],[W0/Wmm W3/Wmm],rp,rs)
%
%[b,a]=ellip(N,rp,rs,Wn,'stop');
%
%dbode(b,a,1/Wm);
%
%\end{verbatim}
%
%\begin{verbatim}
%Wm = 8000;
%t  = (1:500)/Wm;
%s1 = sin(2*pi*t*500);
%s2 = sin(2*pi*t*1500);
%s3 = sin(2*pi*t*2500);
%s = s1 + s2 + s3;
%subplot(5,1,1); plot(t(400:500),s(400:500));
%subplot(5,1,2); plot(t(400:500),s1(400:500));
%subplot(5,1,3); plot(t(400:500),s2(400:500));
%subplot(5,1,4); plot(t(400:500),s3(400:500));
%
%dp = 0.1;
%rs = 40;
%n = 4;
%[b,a]=ellip(4,dp,rs,[1000,2000]*2/Wm);
%sf = filter(b,a,s);
%subplot(5,1,5); plot(t(400:500),sf(400:500));
%\end{verbatim}

